\documentclass[12pt]{article}

\input{../../_assets/preambulo.tex}


\begin{document}

    % 1. Foto de fondo
    % 2. Título
    % 3. Encabezado Izquierdo
    % 4. Color de fondo
    % 5. Coord x del titulo
    % 6. Coord y del titulo
    % 7. Fecha

    
    \input{../../_assets/portada}
    \portadaExamen{ffccA4.jpg}{Ecuaciones\\Diferenciales II\\Examen VIII}{Ecuaciones Diferenciales II. Examen VIII}{MidnightBlue}{-8}{28}{2026}{}

    \begin{description}
        \item[Asignatura] Ecuaciones Diferenciales II.
        \item[Curso Académico] 2021/2022.
        \item[Grado] Doble Grado en Ingeniería Informática y Matemáticas.
        \item[Grupo] Único.
        \item[Profesor] Rafael Ortega Ríos.
        \item[Descripción] Segunda prueba de clase.
        \item[Fecha] 6 de Junio de 2022.
        % \item[Duración] ---.
    
    \end{description}
    \newpage


    % ------------------------------------
    
    \begin{ejercicio}
        Dado un campo vectorial continuo $X : \bb{R} \times \bb{R}^d \to \bb{R}^d$, se supone que $x(t)$ es una solución maximal del problema
        \[
            x' = X(t,x), \quad x(t_0) = x_0,
        \]
        con intervalo maximal $\left]\alpha, \omega\right[$. Se define la función
        \[
            y(t) = x(-t), \quad t \in \left]-\omega, -\alpha\right[.
        \]
        Encuentra un campo continuo $Y : \bb{R} \times \bb{R}^d \to \bb{R}^d$ de manera que $y(t)$ sea solución de
        \[
            y' = Y(t,y), \quad y(-t_0) = x_0.
        \]
        ¿Es $y(t)$ solución maximal? Justifica la respuesta.
    \end{ejercicio}

    \begin{ejercicio}
        Para cada $n \geq 1$ la función $f_n : \left[0, 1\right] \to \bb{R}$ es continua y cumple
        \[
            f_n(t) = \frac{1}{n} + e^t \int_0^t f_n(s)ds.
        \]
        Demuestra que la sucesión $\{f_n\}$ es uniformemente acotada.
    \end{ejercicio}

    \begin{ejercicio}
        La solución del problema de Cauchy
        \[
            x' = x + tx^3, \quad x(0) = x_0
        \]
        se denota por $x(t; x_0)$. Prueba que $x(t; x_0)$ está bien definido en $t \in \left[-1, 1\right]$ cuando $|x_0|$ es suficientemente pequeño. Prueba que, para cada $t \in \left[-1, 1\right]$, se cumple
        \[
            x(t; x_0) = x_0 e^t + o(|x_0|) \quad \text{si } x_0 \to 0.
        \]
    \end{ejercicio}

    \begin{ejercicio}
        Se usa la norma Euclídea en $\bb{R}^3$ y se supone que $X : \bb{R} \times \bb{R}^3 \to \bb{R}^3$ es un campo continuo que cumple
        \[
            \langle X(t,x), x \rangle \leq -2 \quad \text{si } \|x\| = 3.
        \]
        Se supone que $x(t)$ es una solución maximal del problema $x' = X(t,x)$, $x(0) = 0$. Demuestra:
        \begin{enumerate}[label=\alph*)]
            \item $\|x(t)\| < 3$ para todo $t \in \left[0, \omega\right[$
            \item $\omega = +\infty$.
        \end{enumerate}
    \end{ejercicio}

    \begin{ejercicio}
        La función $F : \left[-1, 1\right] \times \bb{R} \to \bb{R}$, $(t,x) \mapsto F(t,x)$ es continua y cumple la condición
        \[
            |F(t, x_1) - F(t, x_2)| \leq L |x_1 - x_2|, \quad \forall (t,x) \in \left[-1, 1\right] \times \bb{R},
        \]
        con $L < \frac{1}{2}$. Para cada $n = 1, 2, \dots$ se supone que $x_n : \left[-1, 1\right] \to \bb{R}$ es una función continua que cumple
        \[
            x_n(t) = \frac{1}{n} + \int_{-t}^t F(s, x_n(s))ds.
        \]
        Demuestra que la sucesión $\{x_n\}_{n \geq 1}$ converge uniformemente en $\left[-1, 1\right]$.
    \end{ejercicio}

    \newpage
    \setcounter{ejercicio}{0} % Reiniciar contador de ejercicios
    % \noindent
    % \textbf{Solución.}

\end{document}